\chapter{Introducción y objetivos}\label{cap:introduccion}

El transporte público metropolitano de Madrid es una infraestructura crítica gestionada
de forma coordinada por el Consorcio Regional de Transportes de Madrid (CRTM), que
integra bajo un mismo marco tarifario y de gobernanza cuatro operadores de naturaleza
muy distinta: Metro de Madrid, la red de superficie de la Empresa Municipal de
Transportes (EMT), las concesiones de autobús interurbano por carretera y los servicios
de Cercanías operados por Renfe. Anticipar con precisión la demanda diaria agregada de
este sistema no es un ejercicio académico: alimenta decisiones operativas reales de
asignación de flota y personal, de dimensionamiento de frecuencias y de planificación
de contingencia ante episodios de demanda anómala (festivos, puentes, eventos
meteorológicos extremos), y lo hace con una urgencia particular en el periodo que cubre
este TFM, íntegramente posterior a la disrupción de movilidad provocada por la pandemia
de COVID-19. Este proyecto adopta, precisamente, una perspectiva que la literatura
revisada en el capítulo \ref{cap:estado_arte} aborda con mucha menor frecuencia de la
que su relevancia práctica exigiría: en lugar de aislar una línea de autobús, una
estación de metro o un operador concreto, se modela la demanda \textbf{sistémica} del
conjunto de los cuatro operadores, preservando en todo momento el desglose por
operador como variable de seguimiento y de ingeniería de características, sobre un
horizonte temporal continuo de 1.310 días (1 de enero de 2023 a 2 de agosto de 2026)
sin un solo hueco ni valor nulo verificado.

La motivación técnica de este TFM es, en última instancia, una pregunta de diseño de
arquitecturas de predicción: ¿mejora la combinación secuencial de un modelo de dinámica
temporal pura y un modelo de corrección exógena sobre residuos a los modelos de una
sola etapa que la literatura ya emplea con éxito en este dominio? La arquitectura
propuesta formaliza esa combinación como un híbrido residual de dos etapas: una red
neuronal recurrente LSTM univariante en la primera etapa, que captura la dinámica
temporal pura de la serie de demanda total, y un modelo XGBoost en la segunda etapa,
que corrige el residuo $e_t = y_t - \hat y^{\text{LSTM}}_t$ a partir de una matriz de
161 variables exógenas de calendario y meteorología retardada, de modo que
$\hat y^{\text{final}} = \hat y^{\text{LSTM}} + \hat e^{\text{XGBoost}}$.

Conviene adelantar aquí, con la misma honestidad con la que se presenta en el capítulo
\ref{cap:resultados}, el resultado central de este trabajo: la arquitectura híbrida,
tal como queda diseñada e implementada, \textbf{no supera} a un modelo SARIMAX
clásico ni a un XGBoost entrenado directamente sobre las mismas variables exógenas
(MAE de test 234.223 frente a 163.627 y 156.121, respectivamente), y el modelo más
preciso de todo el estudio resulta ser un ensamble simple de esos mismos dos componentes
de una sola etapa: 139.725 de MAE de test en su variante equiponderada 50/50 y 139.681 en
su variante ponderada por el inverso del MAE. Este TFM no oculta ese resultado
negativo tras la presentación de su motivación inicial; al contrario, lo convierte en
el objeto de estudio central, auditado con el mismo rigor metodológico con el que se
evaluaría un resultado positivo: validación estrictamente cronológica, protocolo
\emph{out-of-fold} y comparación exhaustiva contra controles.

\section{Objetivo general y específicos}\label{sec:objetivos}

El \textbf{objetivo general} de este TFM es evaluar, bajo un protocolo de validación
estrictamente cronológico y libre de fuga de información, si una arquitectura híbrida
residual LSTM-XGBoost mejora la predicción de la demanda diaria agregada de transporte
público en Madrid frente a baselines estadísticos consolidados y frente a modelos de
aprendizaje automático de una sola etapa entrenados sobre las mismas variables
exógenas, empleando para ello datos reales del CRTM en el periodo 2023-2026.

Este objetivo general se descompone en cinco objetivos específicos, cada uno de los
cuales corresponde a una fase completa del pipeline desarrollado, aunque no todos se
verifican del mismo modo. Esos cinco objetivos consolidan los diez enunciados en la
propuesta inicial del trabajo: las etapas instrumentales que allí figuraban por separado,
el análisis exploratorio y el análisis de importancia de variables, se integran aquí
dentro del objetivo del que son medio y no fin, y se desarrollan en las secciones
\ref{sec:eda} y \ref{sec:importancia_variables} respectivamente; ninguna se ha
suprimido. OE2, OE4 y OE5 son objetivos \emph{comparativos} (enfrentan
modelos entre sí) y quedan respondidos de forma directa y explícita por las
preguntas de investigación PI1-PI5 de la sección \ref{sec:preguntas_investigacion}. OE1
y OE3, en cambio, son objetivos de \textbf{diseño e integridad estructural} del
pipeline: construir una matriz libre de fuga y generar residuos genuinamente fuera de
muestra. No admiten una comparación de desempeño entre modelos y, por tanto, no se
responden mediante una pregunta de investigación, sino que se verifican de forma
exhaustiva mediante la suite de tests y las garantías anti-fuga documentadas en el Anexo
\ref{cap:anexo2}, condición previa y necesaria para que las respuestas a PI1-PI5 sean
fiables:

\begin{itemize}
  \item \textbf{OE1 — Ingestión y matriz de características libre de fuga
    \textnormal{[diseño e integridad; verificado en el Anexo \ref{cap:anexo2}, sin PI
    asociada]}.} Construir un pipeline de ingestión que unifique las tres fuentes
    crudas del proyecto (CRTM, Open-Meteo y el calendario laboral municipal) en una
    serie diaria continua y verificada, y diseñar sobre ella una matriz de 161
    variables predictoras —retardos de demanda total y por operador, medias y
    desviaciones móviles con desplazamiento previo estricto, armónicos de Fourier
    deterministas, codificación exhaustiva de calendario y meteorología retardada—
    sujeta a un conjunto de guardias anti-fuga verificadas por una suite de tests que
    crece con cada fase del proyecto (capítulo \ref{cap:datos_features}).
  \item \textbf{OE2 — Baselines estadísticos y SARIMAX seleccionado por AIC
    \textnormal{[respondido por PI1, PI2]}.}
    Establecer una jerarquía de modelos de referencia —persistencia, persistencia
    estacional, medias móviles y un modelo SARIMAX con exógenas de calendario,
    seleccionado mediante rejilla de criterio de información de Akaike sobre 324
    combinaciones de orden— evaluados bajo un protocolo de pronóstico
    \emph{walk-forward} de un solo paso, de modo que la arquitectura híbrida se mida
    frente a un adversario estadístico ya de por sí exigente y no frente a un baseline
    trivial (sección \ref{sec:baselines}).
  \item \textbf{OE3 — Etapa 1: LSTM univariante con validación \emph{out-of-fold}
    \textnormal{[diseño e integridad; verificado en el Anexo \ref{cap:anexo2}, sin PI
    asociada]}.}
    Entrenar una red LSTM univariante sobre la dinámica temporal pura de la demanda
    total y generar sus predicciones de entrenamiento mediante un esquema
    \emph{walk-forward} expansivo de 5 bloques, de manera que ningún residuo empleado
    en la etapa siguiente proceda de una predicción \emph{in-sample} (sección
    \ref{sec:esquema_oof}), en cumplimiento de una decisión de diseño que el
    documento de contexto técnico del proyecto fija por escrito como no negociable,
    con anterioridad a la obtención de cualquier resultado.
  \item \textbf{OE4 — Etapa 2: XGBoost sobre el residuo \emph{out-of-fold}
    \textnormal{[respondido por PI3, PI4]}.} Entrenar
    un modelo XGBoost sobre los residuos \emph{out-of-fold} sanos de la etapa 1
    utilizando la matriz completa de 161 variables exógenas, y contrastar el resultado
    del híbrido completo frente a dos controles diseñados específicamente para aislar
    si la etapa 1 aporta valor o introduce ruido: la propia etapa 1 en solitario y un
    XGBoost entrenado directamente sobre la demanda total con idéntica matriz de
    variables (secciones \ref{sec:etapa2_xgboost} y \ref{sec:modelos_control}).
  \item \textbf{OE5 — Ensamble, sensibilidad e interpretación teórica del resultado
    \textnormal{[respondido por PI5]}.}
    Explorar, mediante dos experimentos acotados y registrados por adelantado, si un
    repesado de los días de calendario irregular mejora al híbrido y si una
    combinación aritmética simple de los dos modelos de una sola etapa supera a ambos
    por separado, e interpretar el conjunto de resultados a la luz de la teoría de
    combinación de pronósticos de Granger \cite{granger1989combining}, ya citada por
    Zhang \cite{zhang2003hybrid} como condición necesaria para que una hibridación
    mejore a sus componentes (sección \ref{sec:analisis_sensibilidad}).
\end{itemize}

\section{Preguntas de investigación}\label{sec:preguntas_investigacion}

De los tres objetivos específicos de naturaleza comparativa (OE2, OE4 y OE5) se derivan
cinco preguntas de investigación concretas, formuladas aquí con la misma redacción exacta
con la que se responden en la sección \ref{sec:respuesta_preguntas}:

\begin{itemize}
  \item \textbf{PI1.} ¿Bate el híbrido a SARIMAX en validación y en test
    simultáneamente?
  \item \textbf{PI2.} ¿Bate el híbrido a \texttt{xgboost\_alone}, entrenado sobre las
    mismas variables exógenas?
  \item \textbf{PI3.} ¿Reduce el híbrido específicamente el error en festivos y
    puentes frente a SARIMAX y \texttt{xgboost\_alone}?
  \item \textbf{PI4.} ¿Funciona la etapa 2 como mecanismo de corrección residual,
    aunque el híbrido completo no supere a los modelos de una sola etapa?
  \item \textbf{PI5.} ¿Existe una combinación más simple que supere a ambos
    componentes de una sola etapa?
\end{itemize}

\section{Planificación temporal, cronograma y ciclo de vida}\label{sec:planificacion}

El ciclo de vida seguido por este proyecto es, con honestidad, una reconstrucción
documentada del proceso realmente ejecutado y no un cronograma planificado a priori con
fechas de calendario: el documento de contexto técnico del proyecto, versionado y
mantenido a lo largo de todo el desarrollo, no registra fechas absolutas por fase, sino
el cierre
sucesivo de cada una de ellas verificado por una suite de tests que crece de forma
monótona. El patrón de trabajo responde, en esencia, al ciclo clásico de descubrimiento
de conocimiento en bases de datos (KDD) —selección de los datos, integración de las
fuentes, limpieza y preparación, transformación en variables predictoras, modelado,
evaluación e interpretación, y despliegue—, pero aplicado de forma \textbf{incremental y
verificada}: ninguna fase se dio por cerrada
sin que su suite de tests correspondiente pasara en verde, y ninguna fase posterior
reabrió una decisión de diseño de una fase anterior sin dejar constancia explícita del
motivo.

La figura \ref{fig:gantt_fases} y el cuadro \ref{tab:cronograma_fases}
documentan esta secuencia de fases con su eje temporal deliberadamente
relativo (de la Fase 0 a las Fases 7-11) y no en fechas de calendario, en
coherencia con la naturaleza de reconstrucción documentada del proceso y no de
planificación prospectiva. La última etapa de ese ciclo, el \emph{despliegue}, se
materializa en el cuadro de mando interactivo construido en la Fase 6, que reúne los
cinco elementos de visualización comprometidos en la propuesta inicial: demanda real,
demanda predicha, comparación entre modelos, error de predicción e importancia de
variables del modelo XGBoost; el Anexo \ref{cap:anexo3} documenta cómo ejecutarlo.

Tres bloques agrupan las once fases del proyecto. El primero, de preparación de datos
(Fases 0-2), cubre el andamiaje inicial del repositorio, la ingestión y unificación de
las tres fuentes crudas, y la ingeniería de la matriz de 161 variables predictoras. El
segundo, de modelado (Fases 3-5b), concentra el grueso del trabajo experimental de este
TFM: el establecimiento de los baselines estadísticos, el entrenamiento de la etapa 1
LSTM bajo el esquema \emph{out-of-fold}, la construcción y evaluación de la arquitectura
híbrida residual, y los dos experimentos de sensibilidad acotados que cierran la
exploración empírica. El tercero, de explotación y documentación (Fases 6-11), traduce
los artefactos numéricos generados por el bloque anterior en un panel interactivo de
resultados, un refinamiento visual bajo una identidad gráfica coherente, y finalmente
esta misma memoria, modularizada capítulo por capítulo y trazada de principio a fin en
\texttt{docs/\allowbreak{}LaTeX/\allowbreak{}REPORT\_MAPPING.md}.

\begin{figure}[htbp]
  \centering
  \resizebox{\textwidth}{!}{%
  \begin{ganttchart}[
      hgrid,
      vgrid,
      bar/.style={fill=naranja},
      group/.style={fill=gray75},
      x unit=1.0cm,
      % Pasada B (B1): 0.8cm -> 0.5cm. \resizebox{\textwidth}{!} escala por ANCHO, y el
      % ancho natural no depende de y unit; bajar solo la altura de fila acorta el float
      % ~140pt sin tocar el factor de escala, de modo que el cuerpo de letra del diagrama
      % se mantiene intacto (10,61pt medidos) y la figura y el cuadro 1.1 caben en una
      % sola pagina de floats. Reducir en su lugar el ancho del \resizebox habria bajado
      % el cuerpo de letra por debajo del suelo de legibilidad de los incidentes 7/9/10/11.
      y unit chart=0.5cm,
    ]{0}{9}
    \gantttitle{Fase relativa del ciclo de vida}{10} \\
    \gantttitlelist{0,...,9}{1} \\
    \ganttgroup{Preparación de datos}{0}{2} \\
    \ganttbar{Fase 0 -- Andamiaje}{0}{0} \\
    \ganttbar{Fase 1 -- Ingestión}{1}{1} \\
    \ganttbar{Fase 2 -- Ingeniería de características}{2}{2} \\
    \ganttgroup{Modelado}{3}{6} \\
    \ganttbar{Fase 3 -- Baselines}{3}{3} \\
    \ganttbar{Fase 4 -- LSTM etapa 1 (OOF)}{4}{4} \\
    \ganttbar{Fase 5 -- Híbrido residual}{5}{5} \\
    \ganttbar{Fase 5b -- Sensibilidad}{6}{6} \\
    \ganttgroup{Explotación y documentación}{7}{9} \\
    \ganttbar{Fase 6 -- Dashboard}{7}{7} \\
    \ganttbar{Fase 6b -- Refinamiento visual}{8}{8} \\
    \ganttbar{Fases 7-11 -- Redacción de la memoria}{9}{9}
  \end{ganttchart}%
  }
  \caption{\headlinecolor{Ciclo de vida del proyecto por fases (eje relativo: 0 = Andamiaje, 9 = Redacción de la memoria; reconstrucción documentada, sin fechas de calendario).}}
  \label{fig:gantt_fases}
\end{figure}

\begin{table}[htbp]
  \centering
  \footnotesize
  \setlength{\tabcolsep}{4pt}
  \renewcommand{\arraystretch}{0.95}
  \resizebox{\textwidth}{!}{%
    \begin{tabular}{p{1.7cm}>{\raggedright\arraybackslash}p{2.7cm}>{\raggedright\arraybackslash}p{4.3cm}>{\raggedright\arraybackslash}p{4.3cm}r}
\toprule
\rowcolor{naranja}
\textbf{Fase} & \textbf{Denominación} & \textbf{Entregable principal} & \textbf{Artefacto persistido} & \textbf{Pruebas acumuladas} \\
\midrule
\rowcolor{filaclara} Fase 0 & Andamiaje & Estructura src/\allowbreak{}, venv, requirements.txt, contexto técnico del proyecto & Documento de contexto técnico & 0 \\
Fase 1 & Ingesta y unificación & unified\_\allowbreak{}daily.parquet (1.310 x 31) & data/\allowbreak{}interim/\allowbreak{}unified\_\allowbreak{}daily.parquet & 30 \\
\rowcolor{filaclara} Fase 2 & Ingeniería de características anti-fuga & features\_\allowbreak{}daily.parquet (1.310 x 87) & data/\allowbreak{}processed/\allowbreak{}features\_\allowbreak{}daily.parquet & 66 \\
Fase 3 & Particiones, promoción meteorológica y baselines & 105 vars. meteorológicas retardadas, chronological\_\allowbreak{}split, SARIMAX & data/\allowbreak{}processed/\allowbreak{}baseline\_\allowbreak{}predictions.parquet & 111 \\
\rowcolor{filaclara} Fase 4 & LSTM etapa 1 con validación OOF & Walk-forward expansivo 5 folds, lstm\_\allowbreak{}stage1\_\allowbreak{}final.keras & data/\allowbreak{}processed/\allowbreak{}lstm\_\allowbreak{}oof\_\allowbreak{}predictions.parquet & 138 \\
Fase 5 & Híbrido residual & Conjunto residual 552 filas, xgboost\_\allowbreak{}residual.joblib & data/\allowbreak{}processed/\allowbreak{}hybrid\_\allowbreak{}predictions.parquet & 160 \\
\rowcolor{filaclara} Fase 5b & Análisis de sensibilidad prerregistrado & Experimento A (repesado 3x) y B (ensamble) & data/\allowbreak{}processed/\allowbreak{}sensitivity\_\allowbreak{}comparison.parquet & 177 \\
Fase 6 & Cuadro de mando de resultados & Plotly + ipywidgets, 13 PNG & reports/\allowbreak{}figures/\allowbreak{}*.png & 217 \\
\rowcolor{filaclara} Fase 6b & Refinamiento visual & Tema VIU compartido, tablas copiables & src/\allowbreak{}evaluation/\allowbreak{}theme.py & 257 \\
Fases 7-11 & Redacción de la memoria & Modularización LaTeX, figuras y tablas de memoria, REPORT\_\allowbreak{}MAPPING.md & docs/\allowbreak{}LaTeX/\allowbreak{} & 257 \\
\bottomrule
\end{tabular}
%
  }
  \caption{\headlinecolor{Fases del proyecto, entregable principal, artefacto persistido y tests acumulados (fuente: \texttt{export\_latex\_tables.PHASES}).}}
  \label{tab:cronograma_fases}
\end{table}

\begin{sloppypar}
La columna de tests acumulados del cuadro \ref{tab:cronograma_fases} congela su recuento
en el cierre de la Fase 6b (257 pruebas), el último hito de modelado antes de que
comenzara la redacción de la memoria; las Fases 8 y 9 de esa redacción añaden, a su
vez, sus propias suites de verificación (\texttt{test\_memoria\_figures.py} y
\texttt{test\_export\_latex\_tables.py}) que auditan la propia capa de presentación y no
el pipeline de modelado, y que se documentan de forma separada en el anexo
\ref{cap:anexo2}.
\end{sloppypar}

\section{Estructura del documento}\label{sec:estructura_documento}

El resto de la memoria sigue el orden de los objetivos de la sección \ref{sec:objetivos},
en tres bloques. El primero fundamenta el trabajo: los capítulos \ref{cap:estado_arte} y
\ref{cap:datos_features} lo sitúan frente a la literatura y construyen la matriz de
características (OE1). El segundo formaliza el diseño experimental de OE2 a OE5 (capítulo
\ref{cap:metodologia}). El tercero evalúa y concluye (capítulos \ref{cap:resultados} y
\ref{cap:conclusiones}). Cierran la memoria el catálogo de variables (\ref{cap:anexo1}),
la suite de tests (\ref{cap:anexo2}), la reproducibilidad del entorno
(\ref{cap:anexo3}) y la declaración de uso de inteligencia artificial generativa
(\ref{cap:anexo4}).
